\section*{September 7, 2026: Define new construction and review phased groups}
\textit{Agent-drafted entry recording Jacob's decisions and pending evidence.}

Wholly new residential buildings qualify even when they replace demolished
buildings. Jacob confirmed that renovations, additions, and conversions do not
themselves qualify. Parcel subdivision must be traced to identify buildings and
avoid duplicate records; it is not itself construction.

Jacob confirmed combining the two 42-unit buildings at 4400 Grove into one
84-unit observation. The existing commercial decision already takes this
approach, so this confirmation alone changes no measurement or sample.
For Natchez, he requested investigating a combined observation dated to full
completion of the included group. The existing 84-unit record uses 2020, but
neither its full completion year nor its component membership has been newly
verified. A permit issue date and a subsequent completed status do not identify
the date construction finished. Other Natchez records also exist and must not
be merged merely because they share the development's name.

Roosevelt Square remains a grouping question. The recorded review supports one
six-unit building with parcel-specific measurements, but this alone does not
settle whether the intended observation should include other buildings.
Membership, measurement scope, and timing must be established together before
that decision. Jacob requested seeing questionable evidence before consequential
judgments. No general rule merging all phases has been approved.

\small
\textit{Evidence and scope.} These decisions are recorded in the construction
task README. Existing evidence was inspected in
\path{commercial_manual_decisions.csv},
\path{commercial_cross_family_decisions.csv}, and
\path{project_manual_reviews.csv} under that task's adjudication directory.
This entry records definitions and investigation priorities, not a completed
source review. Production code, datasets, and paper estimates are unchanged.
\normalsize
